\section{Results}
\label{sec:results}

All tables in this section are generated directly from the results ledger by
\texttt{scripts/make\_tables.py}; none of the numbers are transcribed. Arms
averaged over fewer than five seeds carry a superscript.

% Generated tables. \IfFileExists keeps the draft building on a machine that
% has the manuscript but not the sweep output.
\IfFileExists{tables/full.tex}{% Generated by scripts/make_tables.py --- do not edit by hand.
% Regenerate after every sweep; the paper reads these directly.

\begin{table}[t]
\centering
\small
\caption{CIFAR-10, 200 epochs, epoch 199. Mean $\pm$ standard deviation over 5 seeds; a superscript marks arms with fewer. KID is from \texttt{torch-fidelity}; see \S\ref{sec:replication} for why the bespoke estimator is not reported.}
\label{tab:full}
\begin{tabular}{lrrrrr}
\toprule
Objective & FID $\downarrow$ & IS $\uparrow$ & KID $\downarrow$ ($\times 10^{3}$) & Prec. $\uparrow$ & Rec. $\uparrow$ \\
\midrule
MSE \emph{(baseline)} & $18.46 \pm 0.21$ & $7.735 \pm 0.070$ & $13.34 \pm 0.27$ & $0.776 \pm 0.002$ & $0.641 \pm 0.003$ \\
$L_1$ & $25.60 \pm 0.17$ & $7.681 \pm 0.057$ & $18.01 \pm 0.20$ & $0.760 \pm 0.003$ & $0.608 \pm 0.004$ \\
Huber & $18.80 \pm 0.02$\,\textsuperscript{2} & $7.769 \pm 0.059$\,\textsuperscript{2} & $13.90 \pm 0.08$\,\textsuperscript{2} & $0.775 \pm 0.001$\,\textsuperscript{2} & $0.647 \pm 0.002$\,\textsuperscript{2} \\
Static wavelet ($p{=}0$) & $18.63 \pm 0.06$ & $7.748 \pm 0.019$ & $13.64 \pm 0.16$ & $0.780 \pm 0.005$ & $0.641 \pm 0.004$ \\
Frequency-aware \emph{(published)} & $18.90 \pm 0.34$ & $7.676 \pm 0.060$ & $13.65 \pm 0.45$ & $0.779 \pm 0.003$ & $0.634 \pm 0.003$ \\
\quad as Eq.~(1) is written & $21.09 \pm 0.29$ & $7.500 \pm 0.023$ & $16.06 \pm 0.24$ & $0.782 \pm 0.006$ & $0.623 \pm 0.004$ \\
\quad normalised to 1 & $19.61 \pm 0.40$ & $7.578 \pm 0.035$ & $14.64 \pm 0.32$ & $0.779 \pm 0.005$ & $0.631 \pm 0.004$ \\
\quad budget held fixed & $19.80 \pm 0.19$ & $7.595 \pm 0.053$ & $14.80 \pm 0.17$ & $0.776 \pm 0.004$ & $0.630 \pm 0.004$ \\
MSE + EMA & $18.46 \pm 0.37$\,\textsuperscript{3} & $7.724 \pm 0.079$\,\textsuperscript{3} & $13.46 \pm 0.25$\,\textsuperscript{3} & $0.778 \pm 0.001$\,\textsuperscript{3} & $0.642 \pm 0.001$\,\textsuperscript{3} \\
Frequency-aware + EMA & $18.95 \pm 0.20$\,\textsuperscript{3} & $7.683 \pm 0.065$\,\textsuperscript{3} & $13.74 \pm 0.34$\,\textsuperscript{3} & $0.776 \pm 0.002$\,\textsuperscript{3} & $0.637 \pm 0.001$\,\textsuperscript{3} \\
\bottomrule
\end{tabular}
\end{table}
}{%
  \begin{table}[t]\centering
  \fbox{\parbox{0.8\linewidth}{\centering\vspace{4mm}
  \textbf{tables/full.tex not built}\\[1mm]
  run \texttt{python scripts/make\_tables.py}\vspace{4mm}}}
  \caption{Full results across objectives and metrics.}\label{tab:full}
  \end{table}}

% The fallback carries the label as well as the box: without it every \ref to
% the table dangles, and a dangling reference is easy to stop noticing.
\IfFileExists{tables/compare_fid.tex}{\input{tables/compare_fid}}{%
  \begin{table}[t]\centering
  \fbox{\parbox{0.8\linewidth}{\centering\vspace{4mm}
  \textbf{tables/compare\_fid.tex not built}\\[1mm]
  run \texttt{python scripts/make\_tables.py}\vspace{4mm}}}
  \caption{Paired comparison against MSE on FID.}\label{tab:compare-fid}
  \end{table}}

Table~\ref{tab:full} reports every objective on every metric;
Table~\ref{tab:compare-fid} gives the paired comparison against MSE. Three
readings matter.

\paragraph{The published objective is not better on any metric.}
Its FID, Inception Score and KID all sit on the wrong side of the baseline, and
none of the three differences clears significance. The point estimates are
consistent with each other in sign, which is worth more than any of them
individually: three metrics computed from two different feature pipelines agree
on the direction.

\paragraph{The ablations that were supposed to isolate the mechanism make it
worse, significantly.} Both normalised variants and the literal reading of the
published equation are significantly worse than MSE after correction, while the
published form itself is not distinguishable from MSE. The ordering is
informative: the closer the objective gets to being \emph{purely} a frequency
weighting, the worse it performs. What keeps the published form competitive is
the timestep reweighting it performs incidentally (\S\ref{sec:confound}).

\paragraph{Seed variance dominates the differences the literature reports.}
The standard deviation of FID across five seeds of an identical configuration
is \num{0.21} for the baseline. Improvements reported in this area are
routinely smaller than that. Any single-run comparison at this scale is
uninformative regardless of which way it happens to fall, which is the
methodological point of \citet{bouthillier2021accounting} instantiated for
diffusion training.

\subsection{Convergence}

The natural fallback for an objective that does not improve final quality is
that it reaches a given quality sooner. It does not: FID is worse at every
evaluated checkpoint (\num{59.9} against \num{52.4} at epoch~39, \num{36.9}
against \num{33.9} at epoch~79, and so on to epoch~199). The gap narrows as
training proceeds but never closes or reverses.

\subsection{Exponential moving averaging}

EMA is reported as its own configuration because it is standard practice and
the published runs omit it, so mixing the two would confound the comparison.
It shifts both arms by less than the seed variance and leaves their ordering
unchanged, which is expected: at a decay of \num{0.9999} over
\num{78200} steps the shadow copy averages roughly the last \num{10000} steps,
by which point the iterate is moving slowly.

\IfFileExists{tables/compare_is.tex}{% Generated by scripts/make_tables.py --- do not edit by hand.
% Regenerate after every sweep; the paper reads these directly.

\begin{table}[t]
\centering
\small
\caption{Paired-by-seed comparison against MSE on Inception Score. $p$ is Holm--Bonferroni corrected across the rows shown. Wilcoxon cannot fall below $0.0625$ at five seeds and is a direction check, not a verdict.}
\label{tab:compare-ismean}
\begin{tabular}{lrrrl}
\toprule
Objective & $\Delta$ & $t$ & $p$ (Holm) & Verdict \\
\midrule
\quad normalised to 1 & $-0.157$ & $-9.98$ & $0.0051$ & \textbf{worse} \\
\quad as Eq.~(1) is written & $-0.235$ & $-7.25$ & $0.0154$ & \textbf{worse} \\
\quad budget held fixed & $-0.140$ & $-5.66$ & $0.0336$ & \textbf{worse} \\
Frequency-aware \emph{(published)} & $-0.059$ & $-1.45$ & $1.0000$ & n.s. \\
Frequency-aware + EMA & $-0.041$ & $-0.62$ & $1.0000$ & n.s. \\
Huber & $-0.009$ & $-0.19$ & $1.0000$ & n.s. \\
$L_1$ & $-0.054$ & $-1.18$ & $1.0000$ & n.s. \\
MSE + EMA & $+0.001$ & $0.08$ & $1.0000$ & n.s. \\
Static wavelet ($p{=}0$) & $+0.013$ & $0.36$ & $1.0000$ & n.s. \\
\bottomrule
\end{tabular}
\end{table}
}{%
  \begin{table}[t]\centering
  \fbox{\parbox{0.8\linewidth}{\centering\vspace{4mm}
  \textbf{tables/compare\_is.tex not built}\\[1mm]
  run \texttt{python scripts/make\_tables.py}\vspace{4mm}}}
  \caption{Paired comparison against MSE on Inception Score.}
  \label{tab:compare-ismean}
  \end{table}}
